\chapter{Nanomaterial Encapsulation, Surface Chemistry \& Multi-Modal Characterization}
\label{chap:nano}

\section{Nanotechnology Taggant Architecture}
While bare synthetic DNA is stable in buffered liquid solutions, harsh industrial environments — such as high-temperature pyrotechnics, metallic projectile casting, explosive detonations, intense UV solar irradiation, and corrosive chemical exposures — demand robust physical protection. To overcome these constraints, the materials science team at SCME-NUST engineered a multi-layered, inorganic core-shell nanoparticle system (Figure~\ref{fig:core_shell_schematic}).

\begin{figure}[H]
    \centering
    % Placeholder for Core-Shell Diagram (Slide 13 / May Report Fig 4.1)
    \fbox{\parbox{0.85\textwidth}{\centering \vspace{2.0cm} \textbf{[Figure 4.1: Multi-Layer Core-Shell Architecture: Silica Core, Cationic CTAB Bilayer, Encapsulated DNA, Phosphor Dopants, and Outer Ceramic Shield]} \vspace{2.0cm}}}
    \caption{Schematic cross-section of the concentric DNA taggant composite nanoparticle.}
    \label{fig:core_shell_schematic}
\end{figure}

The composite taggant particle comprises four concentric functional layers:
\begin{enumerate}
    \item \textbf{Monodisperse Silica Core ($\text{SiO}_2$, 200--300~nm)}: Serves as an isotropic, optically transparent, and chemically inert structural scaffold synthesized via the Stöber sol-gel process.
    \item \textbf{Cationic Surfactant Bilayer (CTAB)}: Reverses the native negative surface charge ($\delta^-$) of silica to an electropositive surface ($\delta^+$), driving spontaneous electrostatic condensation of the polyanionic DNA phosphate backbone.
    \item \textbf{Inorganic Encapsulating Shell ($\text{SiO}_2$)}: Formed by secondary base-catalyzed hydrolysis of tetraethyl orthosilicate (TEOS), permanently entrapping the adsorbed DNA strands within a dense, glassy silica matrix.
    \item \textbf{Luminescent Phosphor Coating \& Outer Ceramic Barrier}: Incorporates rare-earth optical dopants ($\text{SnO}_2:\text{Eu}^{3+}$ and long-afterglow $\text{SrAl}_2\text{O}_4:\text{Eu}^{2+},\text{Dy}^{3+}$) to enable instant UV optical screening, sealed within a refractory ceramic shell withstanding temperatures exceeding $800^\circ\text{C}$.
\end{enumerate}

\section{Chemical Synthesis Methodology}

\subsection{Stage I: Stöber Sol--Gel Synthesis of Monodisperse $\text{SiO}_2$ Nanoparticles}
Spherical silica cores were synthesized using the base-catalyzed Stöber method. Absolute ethanol (99.9\%, 30.0~mL), aqueous ammonia solution (25.0\% $\text{NH}_4\text{OH}$, 3.0~mL), and deionized water (2.0~mL) were homogenized under magnetic stirring at $400\text{ rpm}$ for 10 minutes at ambient temperature ($25^\circ\text{C}$). Tetraethyl orthosilicate (TEOS, 1.5~mL) was introduced in a controlled, dropwise manner over 36 minutes (one drop every 15 seconds). The sol was stirred for 4 hours to ensure quantitative hydrolysis and condensation:
\begin{align}
    \text{Hydrolysis:} \quad & \text{Si}(\text{OC}_2\text{H}_5)_4 + 4\text{H}_2\text{O} \xrightarrow{\text{NH}_3} \text{Si}(\text{OH})_4 + 4\text{C}_2\text{H}_5\text{OH} \\
    \text{Condensation:} \quad & \text{Si--OH} + \text{HO--Si} \longrightarrow \text{Si--O--Si} + \text{H}_2\text{O}
\end{align}
The synthesized nanoparticles were collected by high-speed centrifugation ($8,000\text{ rpm}$, 5 minutes), washed twice in deionized water and twice in ethanol, yielding $0.072\text{ g}$ of pure white $\text{SiO}_2$ nanopowder.

\subsection{Stage II: CTAB Surface Functionalization and Charge Inversion}
Native silica nanoparticles possess a negative zeta potential due to surface deprotonated silanol groups ($\equiv\text{Si--O}^-$). Because DNA is also polyanionic due to its phosphate backbone ($-\text{PO}_4^-$), direct adsorption is prevented by electrostatic repulsion. 

\begin{figure}[H]
    \centering
    % Placeholder for CTAB micelle diagram (Slide 13 / Fig 4.2)
    \fbox{\parbox{0.75\textwidth}{\centering \vspace{1.5cm} \textbf{[Figure 4.2: CTAB Micellar Bilayer Assembly around Silica Core Facilitating Cationic Charge Reversal]} \vspace{1.5cm}}}
    \caption{CTAB self-assembly on silica core providing the cationic ammonium ($\text{N}^+(\text{CH}_3)_3$) interface for DNA binding.}
    \label{fig:ctab_assembly}
\end{figure}

To induce charge reversal, the washed $\text{SiO}_2$ pellet was dispersed in 30.0~mL ethanol by probe ultrasonication (10 minutes). Cetyltrimethylammonium bromide ($\text{CTAB}, \text{C}_{19}\text{H}_{42}\text{BrN}$) was dissolved in ethanol and added dropwise under continuous stirring at $500\text{ rpm}$. The hydrophobic 16-carbon alkyl tails of CTAB pack into a dense bilayer while the positively charged quaternary ammonium head groups ($\delta^+$) orient outward, creating an electropositive interface (Figure~\ref{fig:ctab_assembly}).

\subsection{Stage III: DNA Encapsulation within Secondary Silica Shell}
The CTAB-functionalized $\text{SiO}_2$ suspension was introduced into the buffered synthetic DNA solution at $500\text{ rpm}$, facilitating rapid electrostatic adsorption of the DNA constructs onto the positively charged particle surfaces. Ammonia solution was added to adjust the pH to $\sim 9.0$, activating secondary sol-gel polymerization. 

A secondary volume of TEOS (1.0~mL) was added dropwise over 16 minutes, with an additional 0.6~mL added after 24 hours. The encapsulation reaction proceeded for 2.5 days to allow full growth of the outer amorphous silica network, transforming the suspension into a characteristic high-viscosity protective gel.

\section{Experimental Characterization Results}
The synthesized nanomaterials were subjected to comprehensive structural, chemical, and optical characterization across four complementary modalities: Scanning Electron Microscopy (SEM), Particle Size Analysis (PSA), Fourier Transform Infrared Spectroscopy (FTIR), and X-Ray Diffraction (XRD) (Figure~\ref{fig:four_quadrants}).

\begin{figure}[H]
    \centering
    % Placeholder for 4-Quadrant Characterization (Slide 14)
    \fbox{\parbox{0.9\textwidth}{\centering \vspace{2.2cm} \textbf{[Figure 4.3: Multi-Modal Characterization Quad-Chart: (A) XRD Phase Profiles, (B) PSA Size Distribution, (C) FTIR Spectra, and (D) High-Resolution SEM Micrographs]} \vspace{2.2cm}}}
    \caption{Comprehensive experimental characterization quad-chart of the encapsulated DNA taggants.}
    \label{fig:four_quadrants}
\end{figure}

\subsection{Scanning Electron Microscopy (SEM) and Energy-Dispersive X-Ray Spectroscopy (EDX)}
High-resolution field-emission SEM micrographs (Figure~\ref{fig:four_quadrants}D) obtained at accelerating voltages of $10\text{--}20\text{ kV}$ ($11,000\times$ and $50,000\times$ magnification) revealed highly uniform, spherical nanoparticles with primary diameters tightly distributed between \textbf{200~nm and 300~nm}. EDX spectra confirmed stoichiometric purity, displaying dominant silicon ($K_\alpha = 1.74\text{ keV}$) and oxygen ($K_\alpha = 0.52\text{ keV}$) peaks with zero transition metal contamination.

\subsection{Particle Size Analysis (PSA) by Laser Diffraction}
Laser diffraction particle size analysis (HORIBA LA-920) of the synthesized $\text{SiO}_2$ cores revealed a monodisperse primary distribution centered at $D_{50} = 3.658~\mu\text{m}$ (mean: $3.807~\mu\text{m}$, representing mild controlled agglomeration in suspension). Following DNA adsorption and secondary silica shell growth, a marked rightward hydrodynamic shift toward the \textbf{5--10~$\mu\text{m}$} region was observed (Figure~\ref{fig:four_quadrants}B), providing definitive physical proof of secondary shell deposition encapsulating the DNA-silica complexes.

\subsection{Fourier Transform Infrared Spectroscopy (FTIR)}
FTIR spectra (Figure~\ref{fig:four_quadrants}C) recorded from $4000\text{ cm}^{-1}$ to $400\text{ cm}^{-1}$ conclusively tracked the chemical functionalization at each step:
\begin{itemize}
    \item \textbf{Pure $\text{SiO}_2$}: Exhibited fundamental $\text{Si--O--Si}$ asymmetric stretching at $1080\text{ cm}^{-1}$, $\text{Si--O}$ bending at $465\text{ cm}^{-1}$, and broad surface silanol $\text{O--H}$ stretching at $3400\text{ cm}^{-1}$.
    \item \textbf{CTAB Functionalization}: Introduced characteristic aliphatic $\text{C--H}$ stretching doublets at $2920\text{ cm}^{-1}$ ($\nu_{\text{as}}(\text{CH}_2)$) and $2850\text{ cm}^{-1}$ ($\nu_{\text{s}}(\text{CH}_2)$), confirming surfactant grafting.
    \item \textbf{Post-Encapsulated DNA}: Displayed signature DNA phosphate backbone $\text{P=O}$ asymmetric stretching at $1220\text{--}1240\text{ cm}^{-1}$, nucleobase $\text{C=O}$ vibrations at $1600\text{--}1700\text{ cm}^{-1}$, and deoxyribose $\text{C--O}$ stretching at $1000\text{--}1100\text{ cm}^{-1}$, directly confirming intact chemical preservation of the encapsulated genetic material.
\end{itemize}

\subsection{X-Ray Diffraction (XRD) Phase Analysis}
XRD diffractograms (Figure~\ref{fig:four_quadrants}A) exhibited a classic broad amorphous halo centered at $2\theta \approx 15^\circ\text{--}30^\circ$, confirming the non-crystalline, isotropic network of the Stöber silica matrix. The lack of sharp crystalline peaks confirms that the secondary silica shell grew without localized phase separation, preserving an amorphous cage structure that shields the encapsulated DNA against thermal and mechanical shear.
